\section{Related work}\label{sec:related}

\subsection{Review workflow software}

Rayyan, Covidence, DistillerSR and EPPI-Reviewer share a common design: they
hold the record set, distribute records among reviewers, record include and
exclude decisions, and resolve conflicts
\autocite{Ouzzani_2016,Marshall_2019}. Their contribution is organisational,
and it is substantial: duplicate removal, blind dual screening and conflict
resolution are the mechanisms through which the method protects itself. None
of them, however, produces an automated screening decision that can be
compared with the human one. When a review reports the use of such a tool, the
report describes the workflow, not the performance of an algorithm.

\subsection{Machine learning and active learning}

ASReview introduced an open, actively learned screening loop in which the
reviewer labels records one at a time and the model reorders the remaining
queue after each label \autocite{van_de_Schoot_2021}. SWIFT-Review combined
text mining, clustering and classification in a workbench for scoping
\autocite{Howard_2016}. These systems are evaluated with recall-oriented
metrics, and their reported performance is best understood as the performance
of a prioritisation policy conditioned on the seed labels: change the seeds
and the curve changes. That property is not a defect, but it complicates
transfer, because a review team cannot know in advance what the seed labels
will buy them.

\subsection{Large language models in screening and extraction}

The recent literature moves the decision itself into the model. Khraisha and
colleagues evaluated GPT-4 on screening and data extraction across a set of
reviews and found performance that was usable in some tasks and unreliable in
others \autocite{Khraisha_2024}. Gartlehner and colleagues reached a similar
conclusion for extraction, with a proof-of-concept design
\autocite{Gartlehner_2024}. Reviews of the wider automation landscape stress
that the field lacks shared reporting standards for these evaluations and that
many tools are published without them \autocite{Callaghan_2024}. The journal's
guidance for manuscripts testing generative AI now specifies the minimum:
documented prompts, established metrics, external validation, seeds, model
versions and access dates, and public data and code \autocite{Farotimi_2025}.

\subsection{What is missing}

Three things, in our reading. First, a deployment description: most LLM
screening studies run on a curated benchmark, not inside a review that a team
actually needed to finish. Second, cost accounting at the record level, which
is what determines whether a review of 15{,}000 records is feasible under a
given budget. Third, an instrument that makes the agreement study repeatable:
if computing $\kappa$ between model and reviewer requires a separate script and
a separate dataset, it will be done once, for the paper, and never again. The
platform described next is an attempt to make these three things ordinary.
